\scanpage{7}
\mspage{3}
que $\sigma$ soit l'identité. Pour le voir, on peut supposer que $\sigma$ laisse invariante une structure \deleted{\unsure{de}} de courbe algébrique (de type $(g,\nu)$ anabélien) sur $X$. Il suffit \deleted{\unsure{de}} que l'on ait $\unsure{2g-\nu\geq2|2g+\nu\geq2}$, i.e. que $X$ se plonge dans \unsure{sa} jacobienne généralisée, ext. d'une VA par un tore, pour pouvoir conclure de l'hypothèse homologique que $\sigma$ \deleted{est} induit une \uline{translation} sur cette jacobienne, par un élément d'ordre fini de la jacobienne \unsure{homogène}. Il faut donc voir que $X$ plongé dans sa jacobienne n'est stable par aucune \deleted{\unsure{autom.}} translation non nulle, dans le cas anabélien -- c'est là un fait bien connu\ldots\ [\unsure{Surement|Sûrement}, il ne doit pas être difficile de trouver une démonstration purement topologique de la proposition précédente\ldots]

Notons maintenant que si $Y$ est une surface \inserted{top} anabélienne, tout revêtement étale fini l'est également. On trouve donc le

\deleted{\uline{Proposition}} \inserted{Corollaire} 3\quad Soient $X,Y$ \deleted{deux} comme plus haut, \inserted{$Y$-anabélienne}.
Alors toute classe d'isotopie de morphismes finis étales de $X$ dans $Y$ est un \uline{torseur} à droite sous le groupe $\unsure{\mathcal{A}_X^0|A_X^0}$ des autom de $X$ isotopes à l'identité.

Compte tenu de la proposition 1, il reste
\margin{6}
